% !TeX root = closed-systems-from-comparison-completeness.tex
% ============================================================
\chapter{Electromagnetism from Closedness and Phase Curvature}
\label{ch:em}
% ============================================================
% ============================================================
\begin{chapterorientation}
\orientationitem{Input}{The Hilbert-side scalar channel, connection-first geometry, and the stabilized phase-sensitive quadratic carrier.}
\orientationitem{Role}{This chapter derives the electromagnetic sector as closed phase curvature rather than as an externally supplied field.}
\orientationitem{Output}{The \(U(1)\) phase reduction, Maxwell equations in the realized branch, the charge lattice, and the denominator--\(3\) composite-sector criterion.}
\end{chapterorientation}
\section{Purpose and logical status}
\label{sec:em-purpose}
Under the standing principle of closed-world admissibility
(\cref{sp:closed-world}), this chapter develops the phase-curvature
consequences of the transport-visibility clause
\cref{sp:transport} on the stabilized quadratic carrier.
As the electromagnetic extension of the main argument, it draws on the
transport, quadratic, Hilbert, Einstein-boundary, and connection results
of the preceding chapters to derive phase curvature and charge
quantization.
% ============================================================
The purpose of this chapter is to derive the electromagnetic sector
determined by the closed comparison framework.
The derivation does not proceed in the usual order.
Standard gauge theory typically begins with a field, then introduces a
bundle, then interprets topological classes, and finally extracts
quantization \cite{Nakahara2003}.
The closed-system order is the reverse:
\[
	\begin{aligned}
		 & \text{closedness}
		\Longrightarrow
		\text{finite admissibility}
		\Longrightarrow
		\text{global coherence}
		\\
		 & \Longrightarrow
		\text{discrete phase sectors}
		\Longrightarrow
		\text{field laws}.
	\end{aligned}
\]
\Cref{ch:lift,ch:triangles,ch:interface,ch:hilbert,ch:christoffel}
already established the structural ingredients
needed for this reversal.
\begin{enumerate}
	\item
	      Every admissible enrichment beyond quotient semantics is
	      exhausted by representative choice and morphism-level transport,
	      possibly with both and with no third locus
	      (\cref{cor:two-loci-lift,thm:lift-classification-main}).
	\item
	      The morphism locus is canonically a flat transport connection on
	      protocol categories, while its later global loop-obstruction
	      description is recorded by loop holonomy
	      (\cref{thm:lift-flat-torsor,thm:DO-classification}).
	\item
	      The first nontrivial finite obstruction to transport is ternary:
	      pairwise compatibility does not yet force global compatibility, but a
	      triangle obstruction does (\cref{thm:triangle,thm:bridge-main}).
	\item
	      The first visible stable obstruction layer is the quadratic carrier
	      \[
		      \mathcal K \cong F^2/F^3.
	      \]
	      (\cref{thm:interface-stabilized})
	\item
	      The quadratic carrier admits a canonical orientation involution and a
	      canonical complex structure, unique up to sign
	      (\cref{thm:universal-complex-structure,prop:hilbert-J-unique}).
	\item
	      Under the inherited second-jet faithfulness condition, nonzero
	      stabilized square classes in the degree--\(2\) carrier are
	      equivalent to nonzero realized curvature on the local geometric
	      branch (\cref{thm:interface-curvature}).
	\item
	      Finite comparison patterns are typed at finite Boolean stage, and
	      global existence is equivalent to coherent extension across the
	      refinement tower (\cref{thm:einstein-extension,thm:stability}).
\end{enumerate}
The present chapter combines these facts to prove four statements.
\begin{enumerate}
	\item
	      The morphism locus admits a canonical phase reduction
	      \[
		      \rho_{\mathrm{ph}}:\Omega_{i_0}(\mathcal G(\mathcal I))\to U(1).
	      \]
	\item
	      Its smooth realization determines a phase connection
	      \[
		      A\in\Omega^1(M)
	      \]
	      with curvature
	      \[
		      F=dA\in\Omega^2(M).
	      \]
	\item
	      The electromagnetic field strength \(F=dA\) is the realized
	      phase curvature of the canonical phase reduction, read on the
	      orientation-odd realized degree--\(2\) sector attached to the
	      stabilized quadratic carrier.
	\item
	      Charge quantization is determined by closedness, through finite phase
	      typing and coherent tower extension, before any external topological
	      interpretation is introduced.
\end{enumerate}
These are phase-sector consequences under the chapter hypotheses.  The
claims concern the closed-system phase reduction and its realized
curvature channel; empirical comparison and normalization enter only
through the explicitly marked boundary data and coupling-normalization
clauses below.
Thus electromagnetism is not an additional structure placed on top of
the closed comparison world.
It is the realized phase curvature of the canonical phase reduction of
the morphism locus.
\medskip
The chapter therefore establishes:
\[
	\boxed{
		\begin{aligned}
			\text{Electromagnetism}
				&=
			\text{realized phase curvature}
			\\
				&\quad
			\text{of the canonical phase reduction of the morphism locus}
		\end{aligned}
	}
\]
and
\[
	\boxed{
		\text{charge quantization}=
		\text{closedness of the phase sector}.
	}
\]
% ============================================================
\section{The intrinsic loop representation}
\label{sec:em-loop-representation}
% ============================================================
Let \(\mathcal I\) be a connected protocol category and fix a base object
\(i_0\in\Obj(\mathcal I)\).
\begin{definition}[Loop group]
	The loop group of the transport groupoid of \(\mathcal I\) at \(i_0\) is
	\[
		\Omega_{i_0}(\mathcal G(\mathcal I)):=\Aut_{\mathcal G(\mathcal I)}(i_0).
	\]
\end{definition}
\Cref{thm:lift-classification-main,cor:two-loci-lift} show that every
admissible
distinction beyond quotient-level endpoint data is carried by the
morphism locus, hence by transport.
At the intrinsic level, tree transport is gauge, while loop transport
records the later global loop-obstruction description of the morphism
locus.
\begin{theorem}[Loop-defect representation]
	Under the standing principle \cref{sp:closed-world}, there exists a
	canonical homomorphism
	\[
		\rho_g:\Omega_{i_0}(\mathcal G(\mathcal I))\longrightarrow H_{\mathrm{stab}}
	\]
	into a stabilizer subgroup \(H_{\mathrm{stab}}\) such that:
	\begin{enumerate}
		\item
		      \(\rho_g\) is invariant under gauge modification up to conjugation;
		\item
		      \(\rho_g\) is trivial if and only if transport is endpoint-determined;
		\item
		      all later global loop-obstruction information is carried by
		      \(\rho_g\).
	\end{enumerate}
\end{theorem}
\begin{proof}
	By the classification of representative lifts, every admissible lift
	consists of object representatives together with morphism-level
	transport data.
	Changing representatives modifies the transport cocycle by conjugation,
	so loop holonomy is invariant up to conjugacy.
	If all loop holonomies are trivial, then transport depends only on
	endpoints; conversely, any endpoint-determined transport has trivial
	loop holonomy.
	Thus the later global loop-obstruction description of the morphism
	locus is carried by the loop representation \(\rho_g\).
\end{proof}
% ============================================================
\section{Triangle obstruction and quadratic localization}
\label{sec:em-triangle-localization}
% ============================================================
The loop representation is not yet the electromagnetic sector.
The first task is to identify the precise layer at which intrinsic
transport defect first survives.
\begin{definition}[Detection size]
	Let \(\{C_i\}\) be the family of alignment cosets attached to a lifted
	transport configuration.
	Its detection size is
	\[
		m:=\min\Bigl\{|I|:\ \bigcap_{i\in I} C_i=\varnothing\Bigr\},
	\]
	with \(m=\infty\) if no such finite \(I\) exists.
\end{definition}
\begin{theorem}[Triangle minimality]
	\label{thm:em-triangle-minimal}
	The first nontrivial finite obstruction to simultaneous transport
	compatibility occurs at detection size
	\[
		m=3.
	\]
	Equivalently, pairwise compatibility does not yet force trivial loop
	defect, while a triangle configuration can witness the first genuine
	intrinsic obstruction.
\end{theorem}
\begin{proof}
	By \cref{thm:triangle,thm:bridge-main}, length--\(2\)
	transport relations close trivially, whereas length--\(3\) relations
	produce the first nonvanishing defect.
	Thus no \(2\)-point observable detects a genuinely new obstruction, but
	a \(3\)-point observable can.
	Hence the first finite obstruction is ternary.
\end{proof}
The triangle regime is precisely where the first stable obstruction
carrier appears.
\begin{corollary}[Quadratic localization]
	The intrinsic transport defect factors through the stabilized
	quadratic carrier
	\[
		\mathcal K \cong F^2/F^3.
	\]
\end{corollary}
\begin{proof}
	The triangle obstruction is first visible at degree \(2\), and later
	higher-order terms lie in deeper filtration layers.
	Therefore the first stable intrinsic carrier of visible transport
	defect is \(F^2/F^3\).
\end{proof}
% ============================================================
\section{Canonical phase reduction}
\label{sec:em-phase-reduction}
% ============================================================
The electromagnetic sector must come from the quadratic carrier, not
from an externally imposed gauge group.
The crucial point is that the oriented defect double attached to the
stabilized quadratic carrier already carries a canonical complex
structure.
\begin{theorem}[Canonical phase reduction]
	\label{thm:em-phase-reduction}
	Under the standing principle \cref{sp:closed-world}, there exists a
	canonical abelian phase reduction
	\[
		\rho_{\mathrm{ph}}:\Omega_{i_0}(\mathcal G(\mathcal I))\longrightarrow U(1),
	\]
	unique up to global phase.
\end{theorem}
\begin{proof}
	By \cref{thm:universal-complex-structure}, the oriented quadratic defect
	sector carries a canonical complex structure
	\[
		J:H_{\mathbb R}\to H_{\mathbb R},
		\qquad
		J^2=-\id,
	\]
	unique up to sign.
	This equips the oriented quadratic defect sector with a Hermitian
	structure.
	The abelianization of the loop-defect representation, composed with the
	unitary phase determined by that Hermitian structure, yields a phase
	character
	\[
		\rho_{\mathrm{ph}}:\Omega_{i_0}(\mathcal G(\mathcal I))\to U(1).
	\]
	Since the only ambiguity in the underlying complex structure is a
	global sign, the resulting phase reduction is unique up to global
	phase.
\end{proof}
\begin{remark}
	The appearance of \(U(1)\) is not postulated.
	It is determined by the canonical complex structure on the oriented
	defect double attached to the stabilized quadratic carrier.
\end{remark}
% ============================================================
\section{Orientation involution and the odd transport sector}
\label{sec:em-orientation-odd}
% ============================================================
The phase reduction alone identifies a unitary sector.
To identify electromagnetism itself, one must isolate the
orientation-sensitive part of the quadratic carrier.
\begin{definition}[Loop reversal]
	For a based loop \(\gamma\in\Omega_{i_0}(\mathcal G(\mathcal I))\), define
	\[
		\tau(\gamma):=\gamma^{-1}.
	\]
\end{definition}
\begin{definition}[Orientation involution on the quadratic carrier]
	The loop-reversal map induces a linear involution
	\[
		\tau_{\mathcal K}:\mathcal K\to\mathcal K
	\]
	by functoriality of transport defect under reversal of oriented
	triangle/loop data.
\end{definition}
\begin{lemma}
	\label{lem:em-orientation-square}
	The induced map satisfies
	\[
		\tau_{\mathcal K}^2=\id_{\mathcal K}.
	\]
\end{lemma}
\begin{proof}
	Since \((\gamma^{-1})^{-1}=\gamma\) for every loop \(\gamma\), the
	induced action on any functorially attached defect carrier squares to
	the identity.
\end{proof}
\begin{definition}[Even and odd quadratic sectors]
	Define
	\[
		\mathcal K^+ := \ker(\tau_{\mathcal K}-\id),
		\qquad
		\mathcal K^- := \ker(\tau_{\mathcal K}+\id).
	\]
\end{definition}
Let \(H\) denote the Hilbert completion of the quadratic defect sector
constructed from
\cref{thm:quadratic-scalar,thm:universal-complex-structure}
and its canonical
complex structure.
\begin{theorem}[Canonical orientation decomposition]
	There is a canonical orthogonal decomposition
	\[
		H = H^+ \oplus H^-,
	\]
	where \(H^\pm\) are the closed \(\pm1\)-eigenspaces of the induced
	orientation involution.
\end{theorem}
\begin{proof}
	By \cref{lem:em-orientation-square}, the induced operator on \(H\) is a
	bounded involution.
	Therefore its spectrum lies in \(\{\pm1\}\), and \(H\) decomposes into
	the corresponding eigenspaces
	\[
		H=H^+\oplus H^-.
	\]
	Since the involution is canonical and the Hermitian structure is
	transport-invariant, this decomposition is canonical and orthogonal.
\end{proof}
The odd sector is the one relevant for electromagnetism.
\begin{lemma}[Odd parity determines antisymmetry]
	\label{lem:em-odd-antisymmetry}
	Let \(p\in M\) be a realized point and let \(k\in\mathcal K^-_p\).
	If
	\[
		B_k:T_pM\times T_pM\to\mathbb R
	\]
	is the realized bilinear representative of \(k\), then
	\[
		B_k(v,w)=-B_k(w,v)
		\qquad
		\text{for all }v,w\in T_pM.
	\]
\end{lemma}
\begin{proof}
	Reversal of oriented transport exchanges the ordered pair
	\((v,w)\) with \((w,v)\).
	Since \(k\) lies in the \((-1)\)-eigenspace of the orientation
	involution, the realized degree--\(2\) defect changes sign under this
	exchange.
	Hence
	\[
		B_k(w,v)=-B_k(v,w),
	\]
	so \(B_k\) is alternating.
\end{proof}
\begin{corollary}[Degenerate-pair vanishing]
	\label{cor:em-degenerate-pair-vanishing}
	For every \(k\in\mathcal K^-_p\),
	\[
		B_k(v,v)=0
		\qquad
		\text{for all }v\in T_pM.
	\]
\end{corollary}
\begin{proof}
	Apply \cref{lem:em-odd-antisymmetry} with \(w=v\).
\end{proof}
Thus every odd quadratic class determines an alternating bilinear form.
\begin{definition}[Odd transport \(2\)-form map]
	For each realized point \(p\in M\), define
	\[
		\Xi_p:\mathcal K^-_p\longrightarrow \Lambda^2(T_p^*M)
	\]
	by sending an odd quadratic class to its realized alternating bilinear
	representative.
\end{definition}
\begin{theorem}[Odd sector as \(2\)-form sector]
	\label{thm:em-odd-twoform}
	For each realized point \(p\in M\), the odd quadratic sector admits a
	canonical injective realization
	\[
		\Xi_p:\mathcal K^-_p\hookrightarrow \Lambda^2(T_p^*M).
	\]
	These maps assemble naturally into a bundle morphism
	\[
		\Xi:\mathcal K^-_{\mathrm{real}}\to \Lambda^2(T^*M).
	\]
	If the odd quadratic sector exhausts all realized alternating degree--\(2\)
	transport data, then
	\[
		\mathcal K^-_{\mathrm{real}}\cong \Lambda^2(T^*M).
	\]
\end{theorem}
\begin{proof}
	Under the inherited second-jet faithfulness condition, each odd
	quadratic class on the realized degree--\(2\) channel attached to the
	stabilized quadratic carrier has a realized bilinear representative on
	\(T_pM\times T_pM\), which provides the existence of \(\Xi_p\).
	By \cref{lem:em-odd-antisymmetry} and
	\cref{cor:em-degenerate-pair-vanishing}, the image lies in
	\(\Lambda^2(T_p^*M)\).
	Naturality under the realized local symmetry pseudogroup makes these
	fiberwise maps compatible on overlaps, hence they assemble into a
	bundle morphism.
	Injectivity follows because an odd class with vanishing realized
	alternating representative is invisible in the realized degree--\(2\)
	branch, hence vanishes in the realized odd sector.
	If every realized alternating degree--\(2\) transport datum is produced
	by an odd quadratic class, then the map is also surjective.
\end{proof}
\begin{corollary}[Electromagnetic field strength lives in the odd sector]
	The electromagnetic field strength constructed below is a section of the
	odd transport \(2\)-form bundle.
\end{corollary}
\begin{proof}
	By \cref{thm:em-odd-twoform}, the realized odd quadratic sector embeds
	canonically in \(\Lambda^2(T^*M)\).
	The electromagnetic field strength introduced in
	\cref{def:em-field-strength} is the phase curvature and reverses sign
	under orientation reversal, hence belongs to the odd sector.
\end{proof}
% ============================================================
\section{Phase connection and curvature}
\label{sec:em-connection-curvature}
% ============================================================
The phase reduction determines the electromagnetic connection.
\begin{definition}[Phase line bundle]
	Let \(L_{\mathrm{ph}}\to M\) denote the intrinsic phase line bundle
	determined by the canonical phase reduction.
\end{definition}
\begin{theorem}[Phase connection]
	\label{thm:em-phase-connection}
	Under smooth realization, the phase reduction determines a local unitary
	connection
	\[
		A\in \Omega^1(M).
	\]
\end{theorem}
\begin{proof}
	\Cref{thm:connection-smooth-realization} supplies the
	local transport connection on the realized branch.
	Passing to the phase reduction of the morphism locus yields a unitary
	connection on the intrinsic phase line bundle \(L_{\mathrm{ph}}\).
	In local gauge, this is represented by a \(1\)-form \(A\).
\end{proof}
\begin{definition}[Electromagnetic field strength]
	\label{def:em-field-strength}
	Define
	\[
		F:=dA\in\Omega^2(M).
	\]
\end{definition}
\begin{theorem}[Homogeneous Maxwell equation]
	The electromagnetic field strength satisfies
	\[
		dF=0.
	\]
\end{theorem}
\begin{proof}
	Immediate from \cref{def:em-field-strength}.
\end{proof}
\begin{corollary}
	The field strength \(F\) lies in the odd transport \(2\)-form sector.
\end{corollary}
\begin{proof}
	The phase curvature changes sign under reversal of oriented loop
	transport, hence lies in the orientation-odd part of the degree--\(2\)
	carrier.
	By \cref{thm:em-odd-twoform}, that odd sector is realized as a subbundle
	of \(\Lambda^2(T^*M)\).
\end{proof}
% ============================================================
\section{Admissible phase operators and the sourced Maxwell equation}
\label{sec:em-sourced-maxwell}
% ============================================================
The homogeneous equation is formal.
The sourced equation requires the macroscopic naturality mechanism.
\begin{definition}[Admissible phase operator]
	A phase operator is a map
	\[
		\mathcal E_{\mathrm{ph}}:(g,A)\mapsto J_{\mathrm{ph}}
	\]
	from the realized metric and phase connection to a \(1\)-form
	\(J_{\mathrm{ph}}\),
	satisfying:
	\begin{enumerate}
		\item locality of order at most \(1\) in \(A\);
		\item covariance under the realized local symmetry pseudogroup;
		\item gauge equivariance under
		      \[
			      A\mapsto A+d\chi;
		      \]
		\item dependence only on the intrinsic phase sector of the closed
		      system.
	\end{enumerate}
\end{definition}
\begin{lemma}[Classification of first-order natural phase operators]
	\label{lem:em-natural-operator}
	Let \(F\in\Omega^2(M)\).
	Any linear, first-order, covariant, natural operator that maps \(F\) to
	a \(1\)-form is, up to overall normalization, given by
	\[
		F\mapsto d{*}F.
	\]
\end{lemma}
\begin{proof}
	Naturality permits only constructions from \(g\), \(\nabla\), and \(F\).
	A first-order operator producing a \(1\)-form must contract one index
	of \(\nabla F\) using the metric.
	The only such covariant contraction is
	\[
		\nabla^\mu F_{\mu\nu},
	\]
	which in differential-form notation is exactly \(d{*}F\), up to
	normalization.
\end{proof}
\begin{definition}[Phase current]
	Let \(\psi\) be a charged matter field in a phase representation sector.
	The \emph{phase current} \(J_{\mathrm{ph}}\) is the unique \(1\)-form
	obtained as the natural variation of the matter coupling with respect to
	the phase connection \(A\).
\end{definition}
\begin{theorem}[Sourced Maxwell equation]
	\label{thm:em-sourced-maxwell}
	The electromagnetic field satisfies
	\[
		d{*}F = {*}J_{\mathrm{ph}},
	\]
	where \(J_{\mathrm{ph}}\) is the phase current associated to charged
	matter.
	Equivalently,
	\[
		\nabla_\mu F^{\mu\nu}=g_{\mathrm{ph}}\,J_{\mathrm{ph}}^\nu
	\]
	for a unique coupling normalization \(g_{\mathrm{ph}}\).
\end{theorem}
\begin{proof}
	By \cref{lem:em-natural-operator}, the only admissible first-order
	natural operator on the electromagnetic \(2\)-form is \(d{*}F\), up to
	normalization.
	The matter sector couples to the phase connection and therefore
	produces a covariant phase current \(J_{\mathrm{ph}}\).
	Compatibility of the field equation with locality, covariance, gauge
	equivariance, and first-order dependence determines
	\[
		d{*}F = {*}J_{\mathrm{ph}}
	\]
	up to a single overall constant.
	Writing that constant in index notation gives
	\[
		\nabla_\mu F^{\mu\nu}=g_{\mathrm{ph}}\,J_{\mathrm{ph}}^\nu.
	\]
	No other first-order, gauge-equivariant, covariant law is admissible.
\end{proof}
\begin{corollary}[Charge conservation]
	The phase current is conserved:
	\[
		d{*}J_{\mathrm{ph}}=0,
		\qquad\text{equivalently}\qquad
		\nabla_\nu J_{\mathrm{ph}}^\nu=0.
	\]
\end{corollary}
\begin{proof}
	Apply \(d\) to both sides of \cref{thm:em-sourced-maxwell} and use
	\(d^2=0\).
\end{proof}
% ============================================================
\section{Finite phase typing}
\label{sec:em-finite-phase-typing}
% ============================================================
Up to this point the phase sector is present, but it has not yet been
shown to be discrete.
That discreteness is the real quantization theorem.
\begin{definition}[Finite phase pattern]
	Let \(S\subset \Omega_{i_0}(\mathcal G(\mathcal I))\) be finite.
	A finite phase pattern on \(S\) is a map
	\[
		\vartheta:S\to U(1).
	\]
	It is called \emph{internally admissible} if it is induced by a finite
	comparison configuration of the closed system.
\end{definition}
\begin{definition}[Stage-\(k\) phase sector]
	For each finite Boolean stage \(B_k\) of the refinement tower, let
	\[
		\mathsf P_k
	\]
	denote the set of internally admissible phase sectors represented at
	stage \(k\).
\end{definition}
\begin{theorem}[Finite phase typing theorem]
	\label{thm:em-finite-phase-typing}
	Every internally admissible finite phase pattern induced by the
	canonical phase reduction
	\[
		\rho_{\mathrm{ph}}:\Omega_{i_0}(\mathcal G(\mathcal I))\to U(1)
	\]
	is represented by a finite-stage Boolean comparison pattern.
	Equivalently, for every finite
	\[
		S\subset \Omega_{i_0}(\mathcal G(\mathcal I))
	\]
	and every internally admissible restriction
	\[
		\vartheta=\rho_{\mathrm{ph}}|_S,
	\]
	there exist a finite witness set \(F\subset U\), a finite directed
	comparison pattern \((\pi^\rightarrow,\pi^\leftarrow)\) over \(F\), and
	an index \(k\) such that the realization set
	\[
		X(F;\pi^\rightarrow,\pi^\leftarrow)\in B_k
	\]
	determines \(\vartheta\) on \(S\).
\end{theorem}
\begin{proof}
	We proceed in three steps.
	\smallskip
	\noindent\emph{Step 1: every admissible finite phase restriction is
		transport data.}
	By the two-locus theorem, all admissible non-endpoint distinction lies
	in the morphism locus.
	Hence any admissible finite phase pattern is determined by transport
	data on a finite protocol subconfiguration.
	\smallskip
	\noindent\emph{Step 2: finite transport data is typed by finite
		comparison data.}
	Each loop in a finite set \(S\) is represented by a finite zig-zag of
	morphisms.
	Therefore the transport information relevant to \(S\) depends on only
	finitely many comparison incidences.
	By \cref{thm:einstein-extension}, this finite transport
	information is encoded by a finite directed comparison pattern over a
	finite witness set \(F\subset U\).
	\smallskip
	\noindent\emph{Step 3: finite comparison data lies at finite Boolean
		stage.}
	By \cref{thm:stability}, every finite
	directed comparison pattern belongs to some finite-stage Boolean
	algebra \(B_k\).
	Therefore the phase pattern \(\vartheta\) is represented at finite
	Boolean stage.
	Hence every internally admissible finite phase pattern is finite-stage
	Boolean.
\end{proof}
\begin{corollary}[Discrete finite-stage phase sectors]
	\label{cor:em-discrete-stage-phase}
	For each \(k\), the stage-\(k\) phase sector \(\mathsf P_k\) is
	discrete.
\end{corollary}
\begin{proof}
	Each \(\mathsf P_k\) is represented by admissibility classes inside the
	finite Boolean algebra \(B_k\).
	Hence it is combinatorial rather than continuously parametrized.
\end{proof}
% ============================================================
\section{Closedness-selected quantization}
\label{sec:em-closedness-quantization}
% ============================================================
The phase sector is now finite-stage typed.
Closedness upgrades this to global discreteness.
\begin{theorem}[Closedness-selected phase sectors]
	\label{thm:em-closedness-phase}
	Let
	\[
		\mathsf P_\infty
	\]
	denote the set of globally admissible phase sectors.
	Then restriction along the refinement tower induces a canonical
	identification
	\[
		\mathsf P_\infty \cong \varprojlim_k \mathsf P_k.
	\]
	In particular, the set of globally admissible phase sectors is
	discrete.
\end{theorem}
\begin{proof}
	Global admissibility in the closed system is equivalent to existence of
	a coherent admissible tower.
	By \cref{thm:em-finite-phase-typing}, every admissible finite phase
	restriction lies in some stage-\(k\) sector \(\mathsf P_k\).
	Hence any globally admissible phase sector determines a coherent tower
	\[
		(\vartheta_k)_k,
		\qquad
		\vartheta_k\in\mathsf P_k,
		\qquad
		\vartheta_{k+1}|_{B_k}=\vartheta_k.
	\]
	Conversely, any such coherent tower defines a global admissible phase
	sector by the same inverse-limit mechanism that governs all other
	closed-system objects.
	Therefore
	\[
		\mathsf P_\infty \cong \varprojlim_k \mathsf P_k.
	\]
	Since each \(\mathsf P_k\) is discrete by
	\cref{cor:em-discrete-stage-phase}, the global sector set
	\(\mathsf P_\infty\) is discrete.
\end{proof}
\begin{corollary}[Quantization before geometry]
	Charge quantization is a consequence of closedness before smooth
	geometry is introduced.
	Under smooth realization, the canonical phase reduction realizes
	geometrically through the local phase connection \(A\), whose
	curvature \(F=dA\) exhibits the already-discrete phase sector as a
	quantized phase-curvature sector.
\end{corollary}
\begin{proof}
	By \cref{thm:em-closedness-phase}, the phase sector is already discrete
	at the internal comparison level.
	The canonical phase reduction realizes geometrically through the local
	phase connection \(A\), whose curvature \(F=dA\) exhibits that same
	already-discrete sector on the realized branch.
	Therefore smooth geometry reflects a quantization already determined by
	closedness.
\end{proof}
\begin{remark}[Logical reversal]
	The standard order
	\[
		\text{bundle}
		\Longrightarrow
		\text{topology}
		\Longrightarrow
		\text{quantization}
	\]
	is reversed in the closed-system framework:
	\[
		\begin{aligned}
			 & \text{closedness}
			\Longrightarrow
			\text{finite admissibility}
			\Longrightarrow
			\text{coherent discrete phase sectors}
			\\
			 & \Longrightarrow
			\text{quantized phase curvature}.
		\end{aligned}
	\]
\end{remark}
% ============================================================
\section{Charge lattice and the fundamental spectrum}
\label{sec:em-charge-lattice}
% ============================================================
Discreteness of the phase sector now determines quantization of charges.
\begin{definition}[Global phase sector group]
	Composition of phase sectors induces an abelian group structure on the
	set \(\mathsf P_\infty\) of globally admissible phase sectors.
\end{definition}
\begin{definition}[Primitive phase sector]
	A nonzero phase sector \(\mathfrak p\in \mathsf P_\infty\) is called
	\emph{primitive} if whenever
	\[
		\mathfrak p=\mathfrak a^n
	\]
	for some \(\mathfrak a\in \mathsf P_\infty\) and \(n\in\mathbb Z\), one
	has \(n=\pm1\).
\end{definition}
\begin{definition}[Charge assignment]
	A charge assignment is a group homomorphism
	\[
		q:\mathsf P_\infty\to \mathbb R.
	\]
\end{definition}
\begin{theorem}[Charge lattice theorem]
	\label{thm:em-charge-lattice}
	Assume \(\mathsf P_\infty\) is generated by a primitive phase sector
	\(\mathfrak p_0\).
	Then every admissible phase sector has the form
	\[
		\mathfrak p_n=\mathfrak p_0^{\,n},
		\qquad n\in\mathbb Z,
	\]
	and every charge assignment satisfies
	\[
		q(\mathfrak p_n)=n\,q_0,
		\qquad q_0:=q(\mathfrak p_0).
	\]
	In particular,
	\[
		q(\mathsf P_\infty)=q_0\,\mathbb Z.
	\]
\end{theorem}
\begin{proof}
	Every admissible phase sector is a power of the primitive generator
	\(\mathfrak p_0\).
	Since \(q\) is a group homomorphism,
	\[
		q(\mathfrak p_0^{\,n})=n\,q(\mathfrak p_0)=nq_0.
	\]
	Hence the charge image is exactly the rank--\(1\) lattice
	\(q_0\mathbb Z\).
\end{proof}
\begin{corollary}[Fundamental charge quantization]
	All fundamental charges are integer multiples of a primitive charge unit
	\(q_0\).
\end{corollary}
\begin{proof}
	Immediate from \cref{thm:em-charge-lattice}.
\end{proof}
\begin{remark}
	No fundamental fractional charge arises at the primitive phase level.
	Any effective fractional charge, if present, must therefore be a
	composite-sector quotient effect rather than a fundamental phase
	weight.
\end{remark}
% ============================================================
\section{Composite-sector descent and the minimal denominator}
\label{sec:em-composite-quotients}
% ============================================================
\Cref{sec:em-charge-lattice} closes the fundamental charge lattice.
What remains is the emergent possibility of fractional observable
charge through quotient semantics at the composite level.
\begin{definition}[Composite charged sector]
	A composite charged sector is a tensor product of finitely many
	fundamental charged sectors
	\[
		\mathcal C:=\mathcal S_{n_1}\otimes\cdots\otimes \mathcal S_{n_r},
	\]
	with raw charge
	\[
		q_{\mathrm{raw}}(\mathcal C)=\Bigl(\sum_{j=1}^r n_j\Bigr)q_0.
	\]
\end{definition}
\begin{definition}[Observable descent]
	An observable descent of a composite charged sector is a surjective map
	\[
		\pi_{\mathrm{obs}}:\mathcal C\twoheadrightarrow \mathcal C_{\mathrm{obs}}
	\]
	such that admissible observables on \(\mathcal C\) are precisely those
	that factor through \(\pi_{\mathrm{obs}}\).
\end{definition}
\begin{definition}[Descended charge]
	A descended charge observable on \(\mathcal C_{\mathrm{obs}}\) is a map
	\[
		q_{\mathrm{obs}}:\mathcal C_{\mathrm{obs}}\to\mathbb R
	\]
	such that
	\[
		q_{\mathrm{raw}}=m\,q_{\mathrm{obs}}\circ \pi_{\mathrm{obs}}
	\]
	for some positive integer \(m\).
\end{definition}
\begin{theorem}[Fractional-charge descent criterion]
	\label{thm:em-fractional-criterion}
	A composite observable sector admits an effective fractional charge unit
	\[
		q_{\mathrm{eff}}=\frac{q_0}{m}
		\qquad (m>1)
	\]
	if and only if the raw charge lattice descends through an index-\(m\)
	quotient under observable descent.
\end{theorem}
\begin{proof}
	If
	\[
		q_{\mathrm{raw}}=m\,q_{\mathrm{obs}}\circ\pi_{\mathrm{obs}},
	\]
	then a unit observable charge step corresponds to raw charge
	\(q_0\), hence
	\[
		q_{\mathrm{eff}}=\frac{q_0}{m}.
	\]
	Conversely, if the observable unit is \(q_0/m\), then multiplication by
	\(m\) recovers the raw integer lattice, so the raw charge factors
	through an index-\(m\) quotient.
\end{proof}
The transport spine fixes the minimal possible denominator.
\begin{theorem}[Minimal composite quotient index]
	\label{thm:em-minimal-composite-index}
	If the closed comparison framework determines a nontrivial composite
	observable descent producing effective fractional charge, then the
	minimal possible quotient index is
	\[
		m=3.
	\]
\end{theorem}
\begin{proof}
	By \cref{thm:em-fractional-criterion}, fractional observable charge can
	arise only through a nontrivial quotient of a composite charged sector.
	If \(m=2\), the quotient would already be detectable by pairwise
	compatibility data.
	But the transport spine proves that no genuinely new obstruction is
	visible at arity \(2\); the first irreducible composite obstruction is
	ternary by \cref{thm:em-triangle-minimal}.
	Therefore the first possible determined quotient must occur at index \(3\),
	the first nontrivial composite arity.
\end{proof}
\subsection{Minimal composite descent and uniqueness of denominator 3}
\label{sec:denom3-unique}
\begin{theorem}[Minimal denominator \(3\) and uniqueness]
\label{thm:em-denominator3-uniqueness}
Assume a nontrivial composite observable descent exists.
Then:
\begin{enumerate}[label=\textup{(\roman*)}]
\item the minimal nontrivial denominator is \(3\);
\item any two minimal denominator-\(3\) composite descents are equivalent
      under admissible quotient descent.
\end{enumerate}
\end{theorem}
\begin{proof}
\textbf{Step 1 (transport visibility: \cref{sp:transport}).}
Part \textup{(i)} is exactly
\cref{thm:em-minimal-composite-index} together with
\cref{thm:em-fractional-criterion}.
The first visible nontrivial composite obstruction is ternary, so the
first admissible denominator is forced to be \(3\).

\textbf{Step 2 (intrinsic comparison of minimal descents:
\cref{sp:intrinsic,sp:transport}).}
For part \textup{(ii)}, let
\(\pi_{\mathrm{obs}}^{(1)}\) and \(\pi_{\mathrm{obs}}^{(2)}\)
be two denominator-\(3\) minimal descents.
Minimality forces each to be first-visible at the same ternary composite
level, so their finite comparison tests agree by
\cref{thm:em-triangle-minimal} and the same transport-visibility
criterion used in \cref{thm:em-fractional-criterion}.

\textbf{Step 3 (finite-to-global passage and admissible descent:
\cref{sp:compactness,sp:closure,sp:quotient}).}
By comparison completeness of the closed stack, equality on all finite
comparison tests implies admissible descent equivalence.
Thus the two minimal denominator-\(3\) descents are equivalent.
\end{proof}
\begin{corollary}[Primitive triplet channel]
\label{cor:em-primitive-triplet}
The minimal denominator-\(3\) composite descent determines, up to
admissible equivalence, a primitive irreducible internal multiplicity
channel of dimension \(3\).
\end{corollary}
\begin{proof}
By \cref{thm:em-denominator3-uniqueness}, the first nontrivial composite
descent has unique index \(3\).
Its observable quotient has exactly three primitive residue classes, and
irreducibility follows from minimality (otherwise a proper nontrivial
subchannel would descend at smaller index).
Therefore the primitive channel is three-dimensional and unique up to
admissible equivalence.
\end{proof}
\begin{corollary}[Minimal effective charge unit]
	If a nontrivial composite observable descent is determined, then the minimal
	effective observable charge unit is
	\[
		q_{\mathrm{eff}}=\frac{q_0}{3}.
	\]
\end{corollary}
\begin{proof}
	By \cref{thm:em-fractional-criterion}, effective fractional charge from
	composite descent has unit \(q_0/m\), where \(m\) is the quotient index.
	By \cref{thm:em-minimal-composite-index}, the first admissible index is
	\(m=3\).
	Hence the minimal effective unit is \(q_0/3\).
\end{proof}
\begin{remark}
	Thus thirds are not imported from phenomenology.
	They are the denominator determined by the first irreducible composite
	obstruction size of the transport spine.
\end{remark}
% ============================================================
\section{Uniqueness and normalization of the electromagnetic coupling}
\label{sec:em-normalization}
% ============================================================
The electromagnetic field law is already fixed up to a single coupling
normalization.
We now identify that normalization as internal rather than free.
\begin{definition}[Intrinsic quadratic scale]
	The intrinsic quadratic scale is the normalization of the canonical
	quadratic scalar law
	\[
		Q:\mathcal K\to\mathbb R
	\]
	fixed by the closed-system admissibility structure.
\end{definition}
\begin{definition}[Causal-diamond normalization]
	The causal-diamond normalization is the unique scaling law
	\[
		|D(r)|\sim r^4
	\]
	fixed by the four-dimensional causal-diamond growth theorem.
\end{definition}
\begin{theorem}[Uniqueness of the phase-coupling channel]
	Every admissible abelian phase-curvature law has the form
	\[
		\nabla_\mu \mathcal H^{\mu\nu}=J_{\mathrm{ph}}^\nu,
		\qquad
		\mathcal H^{\mu\nu}=c\,F^{\mu\nu},
	\]
	for a unique nonzero constant \(c\).
	Equivalently, there is exactly one admissible electromagnetic coupling
	channel, up to overall normalization.
\end{theorem}
\begin{proof}
	By \cref{lem:em-natural-operator}, any admissible first-order, natural,
	covariant phase law depends on \(F\) only through \(d{*}F\), up to
	normalization.
	Equivalently, any constitutive tensor must be a scalar multiple of
	\(F\):
	\[
		\mathcal H^{\mu\nu}=cF^{\mu\nu}.
	\]
	Thus there is exactly one coupling channel and one remaining scalar
	normalization.
\end{proof}
\begin{theorem}[Intrinsic phase normalization]
	\label{thm:em-intrinsic-normalization}
	The electromagnetic coupling normalization is determined by the same
	intrinsic normalization mechanism that fixes the quadratic scalar law
	and the causal-diamond scaling.
	In particular, the electromagnetic coupling is not an independent free
	parameter.
\end{theorem}
\begin{proof}
	The first-visible scalar and phase channels inherit their
	normalization from the stabilized quadratic carrier
	\(\mathcal K\cong F^2/F^3\).
	Under the inherited second-jet faithfulness condition, the
	Einstein-side scalar coefficient is read on the realized
	degree--\(2\) channel attached to that stabilized carrier, so the
	scalar, Einstein, and phase channels are normalized relative to one
	shared quadratic scale rather than to independent macroscopic
	inputs.
	Closedness forbids insertion of an additional independent scale.
	The causal-diamond growth law fixes the normalization of macroscopic
	volume and density.
	Therefore the proportionality constant between phase curvature and phase
	current is fixed by the unique conversion between the quadratic carrier
	scale and the causal-diamond density scale.
	Hence the electromagnetic coupling is determined internally.
\end{proof}
\begin{corollary}[Closedness of the electromagnetic sector]
	The electromagnetic field equations are fully determined by the closed
	comparison framework.
	No free parameter remains in the electromagnetic sector.
\end{corollary}
\begin{proof}
	The field strength is fixed by
	\[
		F=dA,
	\]
	the homogeneous law by
	\[
		dF=0,
	\]
	the sourced law by
	\[
		d{*}F={*}J_{\mathrm{ph}},
	\]
	the charge lattice by \cref{thm:em-charge-lattice}, and the coupling
	normalization by \cref{thm:em-intrinsic-normalization}.
	Thus the electromagnetic sector is closed.
\end{proof}
% ============================================================
\section{Structural summary}
\label{sec:em-summary}
% ============================================================
The chapter has proved the following chain:
\[
	\begin{aligned}
		 & \text{morphism locus}
		\Longrightarrow
		\text{loop defect}
		\Longrightarrow
		\text{triangle obstruction}
		\Longrightarrow
		\mathcal K\cong F^2/F^3
		\\
		 & \Longrightarrow
		\text{canonical complex phase}
		\Longrightarrow
		A
		\Longrightarrow
		F=dA,
	\end{aligned}
\]
with
\[
	F\in \Gamma(\Lambda^-_{\mathrm{tr}}(M))\subseteq \Omega^2(M),
\]
and
\[
	\begin{cases}
		dF=0, \\[4pt]
		d{*}F={*}J_{\mathrm{ph}}.
	\end{cases}
\]
At the same time, the internal closedness mechanism yields
\[
	\begin{aligned}
		 & \text{finite phase typing}
		\Longrightarrow
		\text{discrete finite-stage phase sectors}
		\\
		 & \Longrightarrow
		\text{coherent global phase lattice}
		\\
		 & \Longrightarrow
		\text{charge quantization}.
	\end{aligned}
\]
Thus the logic of the chapter is not
\[
	\text{field}
	\Longrightarrow
	\text{bundle}
	\Longrightarrow
	\text{topology}
	\Longrightarrow
	\text{quantization},
\]
but rather
\[
	\begin{aligned}
		 & \text{closedness}
		\Longrightarrow
		\text{finite admissibility}
		\Longrightarrow
		\text{coherent discrete sectors}
		\\
		 & \Longrightarrow
		\text{electromagnetic field laws}.
	\end{aligned}
\]
This is the precise sense in which the closed-system program inverts the
usual order of gauge theory.
\medskip
Accordingly:
\[
	\boxed{
		\begin{aligned}
			&\text{Electromagnetism is the realized phase curvature}\\
			&\text{of the canonical phase reduction of the morphism locus.}
		\end{aligned}
	}
\]
\[
	\boxed{
		\text{Charge quantization is a consequence of closedness.}
	}
\]
\[
	\boxed{
		\text{If effective fractional charge appears, its primitive denominator is }3.
	}
\]
% ============================================================
\section{Uniqueness of the admissible phase operator}
\label{sec:em-operator-uniqueness}
% ============================================================
The sourced Maxwell equation derived in
\cref{sec:em-sourced-maxwell,thm:em-sourced-maxwell} relied on
the classification of admissible phase operators.
We now prove that classification in full tensorial form.
The result is that there is exactly one admissible first-order operator
acting on the electromagnetic field strength.
% ============================================================
\subsection{Statement of the classification theorem}
% ============================================================
\begin{theorem}[Uniqueness of the admissible phase operator]
	\label{thm:em-operator-uniqueness}
	Let $(M,g)$ be the realized spacetime manifold with Lorentzian metric,
	and let
	\[
		F \in \Omega^2(M)
	\]
	be the abelian phase curvature.
	Let
	\[
		\mathcal P:\Omega^2(M)\longrightarrow \Omega^1(M)
	\]
	be a linear operator satisfying:
	\begin{enumerate}
		\item locality of differential order at most $1$;
		\item covariance under the realized local symmetry pseudogroup;
		\item tensorial naturality (no dependence on external structure);
		\item gauge equivariance (dependence only on $F=dA$);
		\item dependence only on the intrinsic phase sector.
	\end{enumerate}
	Then there exists a unique constant $c\in\mathbb R$ such that
	\[
		\mathcal P(F) = c\, d{*}F.
	\]
	Equivalently, in index notation,
	\[
		\mathcal P(F)_\nu = c\, \nabla^\mu F_{\mu\nu}.
	\]
\end{theorem}
% ============================================================
\subsection{Reduction to tensorial classification}
% ============================================================
\begin{lemma}[Gauge reduction]
	\label{lem:em-operator-gauge-reduction}
	Any admissible operator $\mathcal P$ depends on the connection $A$
	only through the curvature $F=dA$.
\end{lemma}
\begin{proof}
	Gauge equivariance requires invariance under
	\[
		A \mapsto A + d\lambda.
	\]
	Since $F=dA$ is gauge-invariant and the theory is abelian, any
	gauge-equivariant expression must factor through $F$.
	Hence $\mathcal P$ depends only on $F$.
\end{proof}
\begin{lemma}[First-order form]
	\label{lem:em-operator-first-order}
	Any admissible operator $\mathcal P$ is a tensorial contraction of
	\[
		\nabla_\alpha F_{\beta\gamma}.
	\]
\end{lemma}
\begin{proof}
	By locality of order at most $1$, $\mathcal P$ may depend on at most
	one derivative of $F$.
	By naturality, the only available structures are $g$, its inverse,
	the Levi-Civita connection, and the volume form.
	Thus $\mathcal P(F)$ must be obtained by contracting
	\[
		\nabla_\alpha F_{\beta\gamma}
	\]
	with these structures.
\end{proof}
% ============================================================
\subsection{Tensorial classification}
% ============================================================
\begin{lemma}[Exhaustion of contractions]
	\label{lem:em-operator-contractions}
	All tensorial contractions of $\nabla_\alpha F_{\beta\gamma}$ producing
	a $1$-form are linear combinations of the following two types:
	\begin{enumerate}
		\item divergence-type contraction:
		      \[
			      \nabla^\mu F_{\mu\nu};
		      \]
		\item antisymmetric $3$-form contraction:
		      \[
			      \nabla_{[\alpha}F_{\beta\gamma]}.
		      \]
	\end{enumerate}
\end{lemma}
\begin{proof}
	The tensor $\nabla_\alpha F_{\beta\gamma}$ has one derivative index and
	two antisymmetric indices.
	To obtain a $1$-form, one must leave exactly one free index and contract
	the remaining two.
	\medskip
	\noindent\emph{Case 1: contract derivative index with one $F$ index.}
	This yields
	\[
		\nabla^\mu F_{\mu\nu}.
	\]
	\medskip
	\noindent\emph{Case 2: antisymmetrize all indices.}
	This yields
	\[
		\nabla_{[\alpha}F_{\beta\gamma]},
	\]
	a $3$-form.
	\medskip
	All other contractions vanish or reduce to these cases:
	\begin{itemize}
		\item contraction of $F$ indices:
		      \[
			      g^{\beta\gamma}F_{\beta\gamma}=0
		      \]
		      by antisymmetry;
		\item contractions using the Levi-Civita tensor reduce to the Hodge
		      dual of the $3$-form above.
	\end{itemize}
	Thus the two listed types exhaust all possibilities.
\end{proof}
% ============================================================
\subsection{Elimination of the antisymmetric term}
% ============================================================
\begin{lemma}[Vanishing of the antisymmetric contraction]
	\label{lem:em-operator-antisymmetric-vanishes}
	The antisymmetric contraction satisfies
	\[
		\nabla_{[\alpha}F_{\beta\gamma]} = 0.
	\]
\end{lemma}
\begin{proof}
	Since $F=dA$, we have
	\[
		dF = d(dA) = 0.
	\]
	In index notation, this is precisely
	\[
		\nabla_{[\alpha}F_{\beta\gamma]} = 0.
	\]
\end{proof}
\begin{corollary}
	The only nonvanishing admissible contraction is
	\[
		\nabla^\mu F_{\mu\nu}.
	\]
\end{corollary}
\begin{proof}
	By \cref{lem:em-operator-contractions}, admissible first-order
	contractions are exhausted by divergence and the fully antisymmetric
	term.
	By \cref{lem:em-operator-antisymmetric-vanishes}, the antisymmetric term
	vanishes identically.
	Therefore only \(\nabla^\mu F_{\mu\nu}\) remains.
\end{proof}
% ============================================================
\subsection{Normalization}
% ============================================================
\begin{lemma}[Constancy of the coefficient]
	\label{lem:em-operator-constant}
	The coefficient multiplying $\nabla^\mu F_{\mu\nu}$ is constant.
\end{lemma}
\begin{proof}
	By naturality, the operator must be invariant under the realized local
	symmetry pseudogroup.
	Hence no coordinate-dependent or field-dependent scalar coefficient is
	admissible.
	The coefficient must therefore be constant.
\end{proof}
% ============================================================
\subsection{Conclusion}
% ============================================================
\begin{proof}[Proof of \cref{thm:em-operator-uniqueness}]
	By \cref{lem:em-operator-gauge-reduction}, $\mathcal P$ depends only on $F$.
	By \cref{lem:em-operator-first-order}, it is built from
	$\nabla F$.
	By \cref{lem:em-operator-contractions}, the only possible contractions
	are divergence-type or antisymmetric.
	By \cref{lem:em-operator-antisymmetric-vanishes}, the antisymmetric term
	vanishes identically.
	Thus
	\[
		\mathcal P(F)_\nu = c\, \nabla^\mu F_{\mu\nu}.
	\]
	By \cref{lem:em-operator-constant}, $c$ is constant.
	In differential-form notation, this is
	\[
		\mathcal P(F)=c\, d{*}F.
	\]
	This proves uniqueness.
\end{proof}
% ============================================================
\subsection{Maxwell equation as a determined law}
% ============================================================
\begin{corollary}[Determined form of the phase field equation]
	The only admissible first-order, local, covariant, gauge-equivariant
	field equation for the electromagnetic $2$-form is
	\[
		d{*}F = {*}J_{\mathrm{ph}},
	\]
	or equivalently
	\[
		\nabla^\mu F_{\mu\nu} = g_{\mathrm{ph}}\, (J_{\mathrm{ph}})_\nu,
	\]
	for a unique coupling constant $g_{\mathrm{ph}}$.
\end{corollary}
\begin{proof}
	By \cref{thm:em-operator-uniqueness}, the left-hand side must be
	proportional to $d{*}F$.
	The right-hand side is the phase current arising from matter coupling.
	Thus the equation follows uniquely up to normalization.
\end{proof}
\begin{remark}[Closure of the electromagnetic sector]
	This classification removes the final ambiguity in the electromagnetic
	sector.
	The field strength $F$, the homogeneous equation $dF=0$, the sourced
	equation $d{*}F={*}J_{\mathrm{ph}}$, and the charge lattice are all
	determined by the
	closed comparison framework.
	No additional admissible first-order phase law exists.
\end{remark}
% ============================================================
\section{Rigidity of the phase symmetry group}
\label{sec:phase-group-rigidity}
% ============================================================
\Cref{thm:em-phase-reduction,thm:em-phase-connection,thm:em-operator-uniqueness}
derive a canonical phase reduction
\[
	\rho_{\mathrm{ph}}:\Omega_{i_0}(\mathcal G(\mathcal I))\to U(1)
\]
from the stabilized quadratic carrier
\[
	\mathcal K \cong F^2/F^3.
\]
The natural question is whether this $U(1)$ phase symmetry is merely one
admissible choice among many, or whether it is the only admissible phase
symmetry determined by the closed comparison stack at the first visible
transport layer.
The answer is rigidity: at degree $2$, the phase symmetry is unique.
% ============================================================
\subsection{Admissible phase symmetry groups}
% ============================================================
\begin{definition}[Admissible phase symmetry group]
	\label{def:admissible-phase-group}
	An \emph{admissible phase symmetry group} at the first visible transport
	layer is a topological group $\Gamma$ together with a homomorphism
	\[
		\rho_\Gamma:\Omega_{i_0}(\mathcal G(\mathcal I))\to \Gamma
	\]
	satisfying:
	\begin{enumerate}
		\item $\rho_\Gamma$ depends only on the stabilized quadratic carrier
		      \[
			      \mathcal K\cong F^2/F^3;
		      \]
		\item $\rho_\Gamma$ is natural under gauge modification and under the
		      realized local symmetry pseudogroup;
		\item $\rho_\Gamma$ is induced by a filtration-compatible structure on
		      the quadratic defect sector;
		\item the corresponding local connection sector is abelian and
		      first-order, so that its curvature law belongs to the admissible phase
		      class of Chapter~\ref{ch:em}.
	\end{enumerate}
\end{definition}
\begin{remark}
	This definition concerns only \emph{phase symmetry} at the first visible
	quadratic layer.
	It does not yet classify all possible symmetry structures that might
	arise on higher transport-jet layers
	\[
		F^m/F^{m+1},\qquad m\ge 3.
	\]
\end{remark}
% ============================================================
\subsection{Rigidity at the quadratic layer}
% ============================================================
\begin{theorem}[Quadratic phase rigidity]
	\label{thm:quadratic-phase-rigidity}
	At the first visible transport layer
	\[
		\mathcal K\cong F^2/F^3,
	\]
	the only admissible phase symmetry group is \(U(1)\), unique up to
	global phase.
	Equivalently, if
	\[
		\rho_\Gamma:\Omega_{i_0}(\mathcal G(\mathcal I))\to \Gamma
	\]
	is an admissible phase symmetry in the sense of
	\cref{def:admissible-phase-group}, then \(\Gamma\) is canonically
	isomorphic to \(U(1)\).
\end{theorem}
\begin{proof}
	The proof is by elimination.
	\smallskip
	\noindent\emph{Step 1: the phase sector must be induced from the quadratic carrier.}
	By hypothesis, an admissible phase symmetry depends only on the first
	visible stabilized transport carrier
	\[
		\mathcal K\cong F^2/F^3.
	\]
	No lower layer can contribute, since degree $1$ carries only first-order
	transport data, and no higher layer can contribute to the first visible
	phase sector without violating the grading of the transport filtration.
	\smallskip
	\noindent\emph{Step 2: the only determined extra structure on the quadratic carrier is the
		canonical complex structure.}
	Chapter~15 proves that the quadratic carrier carries:
	\begin{enumerate}
		\item a unique positive quadratic scalar form;
		\item a unique filtration-compatible orientation involution;
		\item a canonical oriented defect double;
		\item an orthogonal complex structure
		      \[
			      J: H_{\mathbb R}\to H_{\mathbb R},\qquad J^2=-\id,
		      \]
		      unique up to sign;
		\item a unique Hermitian inner product compatible with $J$.
	\end{enumerate}
	Thus the only intrinsic ``phase'' structure determined on the first visible
	carrier is a Hermitian complex line structure, unique up to global sign.
	This is exactly the structure whose unitary symmetry group is \(U(1)\).
	No larger phase group is determined at this layer. \qedhere
\end{proof}
\begin{proof}[Continuation of the proof]
	\smallskip
	\noindent\emph{Step 3: real-linear alternatives are excluded.}
	A merely real-linear phase symmetry would correspond to a structure
	preserving the positive quadratic form without using the canonical
	complex structure.
	But the phase reduction of Chapter~\ref{ch:em} is induced precisely by
	the Hermitian structure determined by \(J\), not just by the underlying real
	quadratic form.
	Hence the phase symmetry must be unitary rather than merely orthogonal.
	In particular, the surviving abelian phase symmetry is not
	\(\mathbb R^\times\), not \(\mathbb R\), and not a generic real torus.
	\smallskip
	\noindent\emph{Step 4: larger abelian phase groups are excluded.}
	Because the complex structure on the quadratic defect sector is unique up
	to sign, there is only one intrinsic phase direction.
	A larger abelian phase group would require several independent phase
	directions on the same first visible layer, equivalently several
	independent compatible complex structures or several independent
	Hermitian phase sectors.
	No such multiplicity is proved or allowed by the quadratic Hilbert
	reconstruction.
	Therefore the abelian phase symmetry at degree $2$ is exactly one copy of
	\(U(1)\).
	\smallskip
	\noindent\emph{Step 5: uniqueness up to global phase.}
	The only residual ambiguity in the complex structure is
	\[
		J\mapsto -J,
	\]
	which corresponds to complex conjugation of the same unitary phase
	sector.
	Therefore the resulting phase reduction is unique up to global phase.
	Hence every admissible phase symmetry group at the first visible
	quadratic layer is canonically isomorphic to \(U(1)\).
\end{proof}
% ============================================================
\subsection{Exclusion of higher-layer phase extensions}
% ============================================================
\begin{theorem}[No independent higher-layer phase extension at degree $2$]
	Let
	\[
		m(F)=F^m/F^{m+1}
	\]
	be the higher transport-jet carriers of Chapter~17.
	Then no admissible phase symmetry at the first visible electromagnetic
	layer can depend on any \(m(F)\) with \(m\ge 3\).
	Equivalently, the electromagnetic phase symmetry is already closed on
	\[
		F^2/F^3.
	\]
\end{theorem}
\begin{proof}
	Chapter~17 proves that for \(m\ge 3\), the graded pieces
	\[
		F^m/F^{m+1}
	\]
	are higher-order transport jets: they encode order--\(m\) Taylor defect
	modulo lower-order data.
	They therefore belong to higher transport structure, not to the first
	visible phase layer.
	An admissible phase symmetry at the electromagnetic level must belong to
	the first visible obstruction layer itself.
	If it depended essentially on some \(F^m/F^{m+1}\) with \(m\ge 3\), then
	the phase sector would not be first-visible at degree \(2\), contrary to
	the construction of the chapter and to the localization
	\[
		\mathcal K\cong F^2/F^3.
	\]
	Hence the electromagnetic phase symmetry is already exhausted on the
	quadratic layer and admits no independent extension from higher
	transport-jet carriers.
\end{proof}
% ============================================================
\subsection{Consequences}
% ============================================================
\begin{corollary}[Uniqueness of electromagnetism at the quadratic layer]
	At the first visible transport layer, electromagnetism is the unique
	admissible phase gauge theory.
\end{corollary}
\begin{proof}
	By \cref{thm:quadratic-phase-rigidity}, the only admissible phase
	symmetry group is \(U(1)\).
	By \cref{thm:em-operator-uniqueness}, the
	unique admissible first-order field law on that \(U(1)\)-sector is the
	Maxwell law
	\[
		dF=0,\qquad d{*}F={*}J_{\mathrm{ph}}.
	\]
	Therefore electromagnetism is the unique admissible phase gauge theory at
	degree \(2\).
\end{proof}
\begin{remark}[What remains open]
	The present theorem does \emph{not} prove that no further non-abelian
	symmetry can arise anywhere in the full transport stack.
	It proves the narrower and sharper claim:
	\[
		\boxed{
			\text{at the first visible quadratic layer, the only admissible phase
				symmetry is }U(1).
		}
	\]
	Any additional gauge structure, if it exists, must therefore arise from
	a genuinely higher transport layer rather than from the electromagnetic
	phase sector itself.
\end{remark}

\section{Conclusion}
\label{sec:em-conclusion}
At the first visible quadratic layer, the electromagnetic conclusion is
rigid: phase symmetry is canonically \(U(1)\), admissible first-order
dynamics is Maxwell, and charge quantization is fixed by the same closure
constraints. No independent degree--\(2\) alternative remains.

This is the forced handoff point of the gauge arc: once the degree--\(2\)
phase sector is closed, any additional internal or matter structure must
come from higher transport forcing rather than from modifying
electromagnetism itself; that continuation is carried by
\cref{ch:matter}.

Accordingly, the scalar channel fixed by the main stack is carried
next through matter-sector classification, numerical closure, and
canonical realization in
\cref{ch:matter,ch:matter-numbers,ch:numerical-closure,ch:matter-realization},
before Appendix \cref{app:fusion} records its downstream
magnetic-confinement application, isolating the canonical intrinsic
transport scalar and separating later perturbative and
confinement-level readings.

\begin{chapterclosing}
\closingitem{Established}{At the first visible quadratic layer, phase symmetry is \(U(1)\), the admissible first-order field law is Maxwell, and the charge lattice includes the denominator--\(3\) composite-sector criterion.}
\closingitem{Carried forward}{Further internal structure must come from higher transport forcing and is treated through matter-sector scalar data.}
\end{chapterclosing}
